% ch1_b (书页 13-24 / p028-p039)
%JOIN%
或
\begin{equation*}
\eta_{\rho\sigma}=\Lambda^{\mu'}{}_{\rho}\,\eta_{\mu'\nu'}\,\Lambda^{\nu'}{}_{\sigma}
=\Lambda^{\mu'}{}_{\rho}\Lambda^{\nu'}{}_{\sigma}\eta_{\mu'\nu'}\,.
\tag{1.29}
\end{equation*}

（在矩阵记号中次序是重要的，而在指标记号中则无关紧要。）我们想要找出矩阵$\Lambda^{\mu'}{}_{\nu}$，使得矩阵$\eta_{\mu'\nu'}$的分量与$\eta_{\rho\sigma}$的分量相同；这就是间隔在这些变换下不变的含义。

满足式(1.28)的那些矩阵称为\textbf{洛伦兹变换}（Lorentz transformations）；它们的全体在矩阵乘法下构成一个群，称为\textbf{洛伦兹群}（Lorentz group）。这个群与SO(3)——三维空间中的转动群——之间有着密切的类比。转动群可以看成由满足$R^{\mathrm{T}}\mathbf{1}R=\mathbf{1}$的$3\times3$矩阵$R$组成，其中$\mathbf{1}$是$3\times3$单位矩阵。这样的矩阵称为\emph{正交}（orthogonal）矩阵，$3\times3$的这种矩阵构成群O(3)。其中不仅有转动，还包括空间坐标轴取向的反转（\textbf{宇称变换}（parity transformations））。有时我们还想排除宇称变换，于是进一步要求矩阵的行列式为1（$|R|=1$）；这样的矩阵称为\emph{特殊}（special）矩阵，所得的群就是SO(3)。如果把它写成

\begin{equation*}
\mathbf{1}=R^{\mathrm{T}}\mathbf{1}R\,,
\tag{1.30}
\end{equation*}

正交条件看上去就更像式(1.28)了。

于是，转动群O(3)与洛伦兹群的差别就在于：把$\mathbf{1}$——所有元素都等于$+1$的$3\times3$对角矩阵——换成$\eta$，一个对角元一个为$-1$、其余为$+1$的$4\times4$对角矩阵。因此洛伦兹群常被称作O(3,1)。它不仅包含推动和转动，还包含时间方向的分立反转以及宇称变换。和前面一样，我们可以要求$|\Lambda|=1$，得到“固有洛伦兹群”（proper Lorentz group）SO(3,1)。然而这并没有给我们真正想要的东西，即连续洛伦兹变换的集合（那些与恒等变换光滑连通的变换），因为时间反转与宇称反转的组合会具有单位行列式。从式(1.29)的$(\rho,\sigma)=(0,0)$分量出发，我们容易证明$|\Lambda^{0'}{}_{0}|\geq1$，其中负值对应时间反转。因此我们最终可以再要求$\Lambda^{0'}{}_{0}\geq1$（加上$|\Lambda|=1$），得到“固有正向”（proper orthochronous）或“受限”（restricted）洛伦兹群。有时它被记作类似$SO(3,1)\uparrow$的符号，但通常我们不会费心去明确作出这种区分。注意，$3\times3$单位矩阵其实就是普通平直空间的度规。所有本征值都为正的这种度规称为\textbf{欧几里得型}（Euclidean）度规，而像式(1.15)那样只带一个负号的度规则称为\textbf{洛伦兹型}（Lorentzian）度规。

写下简单洛伦兹变换的显式表达式是直截了当的。$x$-$y$平面中一个我们熟悉的转动是

\begin{equation*}
\Lambda^{\mu'}{}_{\nu}=
\begin{pmatrix}
1 & 0 & 0 & 0\\
0 & \cos\theta & \sin\theta & 0\\
0 & -\sin\theta & \cos\theta & 0\\
0 & 0 & 0 & 1
\end{pmatrix}.
\tag{1.31}
\end{equation*}

转动角$\theta$是周期为$2\pi$的周期变量。推动则可以想成“空间方向与时间方向之间的转动”。沿$x$方向的推动就是一个例子：

\begin{equation*}
\Lambda^{\mu'}{}_{\nu}=
\begin{pmatrix}
\cosh\phi & -\sinh\phi & 0 & 0\\
-\sinh\phi & \cosh\phi & 0 & 0\\
0 & 0 & 1 & 0\\
0 & 0 & 0 & 1
\end{pmatrix}.
\tag{1.32}
\end{equation*}

推动参量$\phi$与转动角不同，其定义范围是从$-\infty$到$\infty$。把各个变换相乘就可以得到一般的变换；这个六参数矩阵（三个推动、三个转动）的显式表达式并不好看，也不值得费神写下来。一般说来洛伦兹变换不可对易，所以洛伦兹群是非阿贝尔群（nonabelian）。平移与洛伦兹变换合在一起构成一个十参数非阿贝尔群，即\textbf{庞加莱群}（Poincaré group）。

大家得知推动对应于通过换到一个以恒定速度运动的参考系来改变坐标时，不应感到惊讶，但我们不妨更明白地看一下。（不要把“推动”（boosting）与“加速”（accelerating）混淆。变到一个不同参考系的推动与加速一个物体这两者之间的差别，就如同转到不同坐标系与让物体自旋这两者之间的差别。）对于由式(1.32)给出的变换，变换后的坐标$t'$和$x'$将为

\begin{align*}
t'&=t\cosh\phi-x\sinh\phi\\
x'&=-t\sinh\phi+x\cosh\phi\,.
\tag{1.33}
\end{align*}

由此我们看到，由$x'=0$定义的那个点在运动；它具有速度

\begin{equation*}
v=\frac{x}{t}=\frac{\sinh\phi}{\cosh\phi}=\tanh\phi\,.
\tag{1.34}
\end{equation*}

为了换成更平实的记号，我们可以代入$\phi=\tanh^{-1}v$，得到

\begin{align*}
t'&=\gamma(t-vx)\\
x'&=\gamma(x-vt)\,,
\tag{1.35}
\end{align*}

其中$\gamma=1/\sqrt{1-v^{2}}$。的确如此，我们这条抽象路线重新得到了洛伦兹变换的惯用表达式。应用这些公式就得到时间膨胀（time dilation）、长度收缩（length contraction）等等。

在时空图（spacetime diagram）的语境中考察洛伦兹变换是很有启发性的。根据式(1.33)，在$x$-$t$平面内的推动下，$x'$轴（$t'=0$）由$t=x\tanh\phi$给出，而$t'$轴（$x'=0$）则由$t=x/\tanh\phi$给出。于是我们看到，空间轴和时间轴彼此向对方转过去，只不过它们像剪刀一样斜着合拢，而不再保持传统欧几里得意义下的正交。（我们将会看到，这些轴事实上在洛伦兹意义下确实仍保持正交；这正是度规在推动下保持不变的含义。）

%FIG{1.7}: {洛伦兹变换把$\{t',x'\}$坐标与$\{t,x\}$坐标联系起来。注意光锥不变。}

这应该毫不奇怪，因为假如时空的行为就像四维版本的空间，那世界将会是一个非常不同的地方。我们相当鲜明地看到了这种情形与Newton世界之间的区别；在狭义相对论中，我们不可能（以不依赖于坐标的方式）说出一个与点$p$类空间隔的点究竟是在$p$的未来、$p$的过去，还是“与$p$同时”。

还要注意，由$x'=\pm t'$定义的路径与由$x=\pm t$定义的路径完全相同；这些轨迹在沿$x$轴的推动下保持不变。当然，我们知道光正是以这个速度传播的；于是我们便发现，光速在任何惯性系中都是相同的。

\section{矢量}

要更细致地探究 Minkowski 空间的结构，就需要引入矢量与张量的概念。我们先从矢量讲起，大家应该对它已经熟悉。当然，时空中的矢量是四维的，常被称为\textbf{四维矢量}（four-vectors）。这会带来相当可观的差别——例如，两个四维矢量之间根本不存在叉积这种东西。

除了维数这个简单事实之外，最需要强调的一点是：每个矢量都定域在时空中的一个给定点上。你可能习惯于把矢量看成从空间中一点拉伸到另一点的东西，甚至看成可以随意从一个点滑动到另一个点的“自由”矢量。在平直空间的语境之外，这些都不是有用的概念；一旦引入曲率，我们就失去了从一个点向另一个点画出特殊曲线的能力，也失去了把矢量唯一地搬动绕行流形的能力。毋宁说，对时空中的每个点$p$，我们都把定域于该点的所有可能矢量的集合与它联系起来；这个集合称为$p$处的\textbf{切空间}（tangent space），记作$T_p$。这个名称的灵感来自这样的图像：在一个简单的弯曲二维空间中，附着于某点的那些矢量全体构成一个在该点与此空间相切的平面。（这幅图像依赖于把流形和切空间嵌入一个更高维的外部空间，而一般来说我们既没有也不需要这种嵌入。）抛开这个灵感不谈，重要的是把这些矢量看成定域于单独一点，而不是从一点拉伸到另一点（尽管这不会妨碍我们在时空图上把它们画成箭头）。

在第2章中，我们将把每个点处的切空间与我们可以从时空本身构造出来的东西联系起来。眼下，只需把$T_p$想成时空每一点上的一个抽象矢量空间。所谓\textbf{（实）矢量空间}（real vector space），是指由一些对象（矢量）组成的集合，它们可以线性地彼此相加并乘以实数。于是，对任意两个矢量$V$、$W$和实数$a$、$b$，我们有

\begin{equation*}
(a+b)(V+W)=aV+bV+aW+bW\,.
\tag{1.36}
\end{equation*}

每个矢量空间都有一个原点，即在矢量加法下充当恒等元的零矢量。在许多矢量空间中还有别的运算，比如取内积（点积），但这是矢量空间这一初等概念之外的额外结构。

矢量是一个定义得非常明确的几何对象，\textbf{矢量场}（vector field）也一样，它定义为在时空每一点上恰好有一个矢量的矢量的集合。[把一个$n$维流形$M$的所有切空间汇集起来，可以组装成一个$2n$维流形，称为\textbf{切丛}（tangent bundle），$T(M)$。它是“纤维丛”（fiber bundle）的一个具体例子，纤维丛还带有某种额外的数学结构；就目前的目的我们不需要那些细节。]尽管如此，把矢量相对于某组基矢分解成分量往往是有用的。\textbf{基}（basis）是任何一组既张满矢量空间（任何矢量都是基矢的线性组合）又线性无关（基中没有任何一个矢量是其余基矢的线性组合）的矢量。对任何给定的矢量空间，可供选择的基有无穷多种，但每个基都由相同数目的矢量组成，这个数目称为该空间的\textbf{维数}（dimension）。（对于 Minkowski 空间中一点相应的切空间，维数当然是四。）

%FIG{1.8}: {切空间$T_p$的示意性图示——点$p$处所有矢量构成的空间。}

设想在每个切空间中都立起一组四个矢量$\hat{e}_{(\mu)}$作为基，$\mu$和通常一样取遍$\{0,1,2,3\}$。实际上，不妨说每个基都是“适配于坐标$x^\mu$”的——也就是说，基矢$\hat{e}_{(1)}$就是我们通常认为指向$x$轴方向的那个矢量。我们完全没有必要选择适配于任何坐标系的基，尽管这样做往往很方便。（和前面一样，我们在这里本可以更精确一些，但后面我们会以令人难以忍受的精确程度重新讨论这个问题，所以现在稍微马虎一点是可以原谅的。）于是，任何抽象矢量$A$都可以写成基矢的线性组合：

\begin{equation*}
A=A^{\mu}\hat{e}_{(\mu)}\,.
\tag{1.37}
\end{equation*}

系数$A^{\mu}$就是矢量$A$的\textbf{分量}（components）。我们多半会把基完全忘掉，而有点随意地把“矢量$A^{\mu}$”挂在嘴边，但要记住这只是一种简写。真正的矢量是抽象的几何实体，而分量只是矢量在某个方便的基下基矢的系数。（由于我们通常会不写出显式基矢，指标一般就用来标记矢量和张量的分量。这就是基矢的指标要加括号的原因——提醒我们这是矢量的集合，而不是单个矢量的分量。）

时空中的一个矢量的标准例子是曲线的切矢量。穿过时空的一条参数化曲线或路径由坐标作为参量的函数给出，例如$x^{\mu}(\lambda)$。切矢量$V(\lambda)$的分量为

\begin{equation*}
V^{\mu}=\frac{\mathrm{d}x^{\mu}}{\mathrm{d}\lambda}\,.
\tag{1.38}
\end{equation*}

完整的矢量是$V=V^{\mu}\hat{e}_{(\mu)}$。在洛伦兹变换下，坐标$x^{\mu}$按照式(1.25)变换，而参量化$\lambda$不变；因此我们可以推出，切矢量的分量必须按如下方式变化

\begin{equation*}
\boxed{\,V^{\mu}\rightarrow V^{\mu'}=\Lambda^{\mu'}{}_{\nu}V^{\nu}\,.}
\tag{1.39}
\end{equation*}

然而，矢量$V$本身（与其在某个坐标系中的分量相对而言）在洛伦兹变换下是不变的。我们可以利用这个事实来导出基矢的变换性质。把变换后坐标系中的那组基矢记为$\hat{e}_{(\nu')}$。由于矢量不变，我们有

\begin{equation*}
V=V^{\mu}\hat{e}_{(\mu)}=V^{\nu'}\hat{e}_{(\nu')}
=\Lambda^{\nu'}{}_{\mu}V^{\mu}\hat{e}_{(\nu')}\,.
\tag{1.40}
\end{equation*}

但无论分量$V^{\mu}$的数值是什么，这个关系都必须成立。因此我们可以说

\begin{equation*}
\hat{e}_{(\mu)}=\Lambda^{\nu'}{}_{\mu}\hat{e}_{(\nu')}\,.
\tag{1.41}
\end{equation*}

要把新基$\hat{e}_{(\nu')}$用旧基$\hat{e}_{(\mu)}$表示出来，我们应该乘以洛伦兹变换$\Lambda^{\nu'}{}_{\mu}$的逆。但是，从无撇坐标变到带撇坐标的洛伦兹变换的逆也是一个洛伦兹变换，只不过这次是从带撇系变到无撇系。因此我们将引入一种略显精妙的记号：对这两个矩阵用同一个符号，只是把带撇指标与无撇指标互换。也就是说，由$\Lambda^{\mu'}{}_{\nu}$规定的洛伦兹变换，其逆变换写作$\Lambda^{\rho}{}_{\sigma'}$。从运算上讲，这意味着

\begin{equation*}
\Lambda^{\mu}{}_{\nu'}\Lambda^{\nu'}{}_{\rho}=\delta^{\mu}{}_{\rho}\,,\qquad
\Lambda^{\sigma'}{}_{\lambda}\Lambda^{\lambda}{}_{\tau'}=\delta^{\sigma'}{}_{\tau'}\,.
\tag{1.42}
\end{equation*}

于是由式(1.41)我们得到基矢的变换规则：

\begin{equation*}
\hat{e}_{(\nu')}=\Lambda^{\mu}{}_{\nu'}\hat{e}_{(\mu)}\,.
\tag{1.43}
\end{equation*}

因此，基矢的集合按坐标或矢量分量的洛伦兹变换之逆来变换。

让我们稍停片刻，把这一切消化一下。我们引入了用上指标标记的坐标，它们在洛伦兹变换下按某种方式变换。接着我们考虑了也用上指标写出的矢量分量，这是有意义的，因为它们以与坐标函数相同的方式变换。（在固定的坐标系中，四个坐标$x^{\mu}$中的每一个都可以看成时空上的函数，矢量场的四个分量中的每一个也是如此。）与坐标系相联系的基矢则通过逆矩阵变换，并用下指标标记。这种记号保证了由分量与基矢求和构造出来的不变对象在变换下保持不变，正如我们所愿。如果我说，对于可能有多个指标的张量，情形也将继续如此，大概不算把后面的内容剧透太多。

\section{对偶矢量（1-形式）}

一旦建立了一个矢量空间，我们就可以定义另一个与之相联系的矢量空间（维数相同），称为\textbf{对偶矢量空间}（dual vector space）。对偶空间通常用星号标记，于是切空间$T_p$的对偶空间称为\textbf{余切空间}（cotangent space），记作$T_p^{*}$。对偶空间是从原来的矢量空间到实数的所有线性映射构成的空间；用数学行话说，若$\omega\in T_p^{*}$是一个对偶矢量（dual vector），则它作为一个映射来作用，满足

\begin{equation*}
\omega(aV+bW)=a\omega(V)+b\omega(W)\in\mathbf{R}\,,
\tag{1.44}
\end{equation*}

其中$V$、$W$是矢量，$a$、$b$是实数。这些映射的妙处在于它们本身也构成一个矢量空间；于是，若$\omega$和$\eta$是对偶矢量，我们有

\begin{equation*}
(a\omega+b\eta)(V)=a\omega(V)+b\eta(V)\,.
\tag{1.45}
\end{equation*}

为了让这个构造更具体一点，我们可以引入一组基对偶矢量$\hat{\theta}^{(\nu)}$，办法是要求

\begin{equation*}
\hat{\theta}^{(\nu)}(\hat{e}_{(\mu)})=\delta^{\nu}{}_{\mu}\,.
\tag{1.46}
\end{equation*}

于是每个对偶矢量都可以用它的分量写出，我们用下指标来标记分量：

\begin{equation*}
\omega=\omega_{\mu}\hat{\theta}^{(\mu)}\,.
\tag{1.47}
\end{equation*}

通常，与矢量的情形完全类似，我们会径直写$\omega_{\mu}$，用它代表整个对偶矢量。实际上，你有时会看到$T_p$的元素（就是我们所说的矢量）被称为\textbf{逆变矢量}（contravariant vectors），而$T_p^{*}$的元素（就是我们所说的对偶矢量）被称为\textbf{协变矢量}（covariant vectors），不过在如今这个时代，这些术语听起来有点过时了。只要你把普通矢量称作带上指标的矢量、把对偶矢量称作带下指标的矢量，应该没有人会不高兴。对偶矢量的另一个名字是\textbf{1-形式}（one-forms），这个多少有些神秘的称呼在第2章中会变得更清楚。

分量记号给出了一种书写对偶矢量作用于矢量的简单方式：

\begin{align*}
\omega(V)&=\omega_{\mu}\hat{\theta}^{(\mu)}(V^{\nu}\hat{e}_{(\nu)})\\
&=\omega_{\mu}V^{\nu}\hat{\theta}^{(\mu)}(\hat{e}_{(\nu)})\\
&=\omega_{\mu}V^{\nu}\delta^{\mu}{}_{\nu}\\
&=\omega_{\mu}V^{\mu}\in\mathbf{R}\,.
\tag{1.48}
\end{align*}

这正是我们很少需要把基矢和对偶矢量显式写出来的原因；分量包办了一切工作。式(1.48)的形式还暗示，通过定义

\begin{equation*}
V(\omega)\equiv\omega(V)=\omega_{\mu}V^{\mu}\,.
\tag{1.49}
\end{equation*}

我们可以把矢量看成对偶矢量上的线性映射。因此，对偶矢量空间的对偶空间就是原来的矢量空间本身。

当然，在时空中我们感兴趣的不是一个单独的矢量空间，而是矢量场和对偶矢量场。[把$M$上所有余切空间汇合起来，就得到\textbf{余切丛}（cotangent bundle），$T^{*}(M)$。]在那种情形下，对偶矢量场作用在矢量场上得到的不是一个单独的数，而是时空上的一个\textbf{标量}（scalar）（函数）。标量是不带指标的量，它在洛伦兹变换下不变；它是从时空到实数的、不依赖于坐标的映射。

我们可以用前面处理矢量时用过的同样论证（几何对象独立于坐标，即使它们的分量并非如此）来导出对偶矢量的变换性质。答案是：对于分量，

\begin{equation*}
\boxed{\,\omega_{\mu'}=\Lambda^{\nu}{}_{\mu'}\omega_{\nu}\,,}
\tag{1.50}
\end{equation*}

而对于基对偶矢量，

\begin{equation*}
\hat{\theta}^{(\rho')}=\Lambda^{\rho'}{}_{\sigma}\hat{\theta}^{(\sigma)}\,.
\tag{1.51}
\end{equation*}

这正是从指标位置所能预期的；对偶矢量的分量按矢量分量的逆变换来变换。注意这就保证了标量(1.48)在洛伦兹变换下不变，正如它本应如此。

在时空中，对偶矢量最简单的例子是标量函数的\textbf{梯度}（gradient），即对时空坐标的偏导数的集合，我们用小写的d来记它：

\begin{equation*}
\mathrm{d}\phi=\frac{\partial\phi}{\partial x^{\mu}}\hat{\theta}^{(\mu)}\,.
\tag{1.52}
\end{equation*}

用来变换偏导数的常规链式法则，在这种情形下恰好就是对偶矢量分量的变换法则：

\begin{align*}
\frac{\partial\phi}{\partial x^{\mu'}}&=\frac{\partial x^{\mu}}{\partial x^{\mu'}}\frac{\partial\phi}{\partial x^{\mu}}\\
&=\Lambda^{\mu}{}_{\mu'}\frac{\partial\phi}{\partial x^{\mu}}\,,
\tag{1.53}
\end{align*}

其中我们用了式(1.25)把洛伦兹变换与坐标联系起来。梯度是对偶矢量这一事实，导致了偏导数的如下简写记号：

\begin{equation*}
\frac{\partial\phi}{\partial x^{\mu}}=\partial_{\mu}\phi=\phi_{,\mu}\,.
\tag{1.54}
\end{equation*}

于是，$x^{\mu}$带上指标，但当它出现在导数的分母中时，它意味着所得的对象带下指标。在本书中我们一般使用$\partial_{\mu}$而不用逗号记号。注意，对于前面给出的那个矢量例子——曲线的切矢量，梯度确实以一种自然的方式作用。其结果是函数沿曲线的普通导数：

\begin{equation*}
\partial_{\mu}\phi\,\frac{\partial x^{\mu}}{\partial\lambda}=\frac{\mathrm{d}\phi}{\mathrm{d}\lambda}\,.
\tag{1.55}
\end{equation*}

\section{张量}

矢量与对偶矢量的一个直接推广就是\textbf{张量}（tensor）的概念。正如对偶矢量是从矢量到$\mathbf{R}$的线性映射，一个型（或阶）$(k,l)$的张量$T$是从一组对偶矢量和矢量到$\mathbf{R}$的多重线性映射（multilinear map）：

\begin{equation*}
T:\;T_p^{*}\times\underset{(k\text{ 次})}{\cdots}\times T_p^{*}
\times T_p\times\underset{(l\text{ 次})}{\cdots}\times T_p\rightarrow\mathbf{R}\,.
\tag{1.56}
\end{equation*}

这里，“$\times$”表示笛卡尔积（Cartesian product），例如$T_p\times T_p$就是有序矢量对构成的空间。多重线性意味着张量作用在它的每一个宗量上都是线性的；例如，对于型$(1,1)$的张量，我们有

\begin{align*}
T(a\omega+b\eta,\,cV+dW)={}&acT(\omega,V)\\
&+adT(\omega,W)+bcT(\eta,V)+bdT(\eta,W)\,.
\tag{1.57}
\end{align*}

从这个观点来看，标量是型$(0,0)$的张量，矢量是型$(1,0)$的张量，而对偶矢量是型$(0,1)$的张量。

固定型$(k,l)$的所有张量构成一个矢量空间；它们可以彼此相加并乘以实数。为了给这个空间构造基，我们需要定义一种新的运算，称为\textbf{张量积}（tensor product），记作$\otimes$。若$T$是$(k,l)$张量，$S$是$(m,n)$张量，我们如下定义$(k+m,\,l+n)$张量$T\otimes S$：

\begin{align*}
&T\otimes S(\omega^{(1)},\ldots,\omega^{(k)},\ldots,\omega^{(k+m)},V^{(1)},\ldots,V^{(l)},\ldots,V^{(l+n)})\\
&\qquad =T(\omega^{(1)},\ldots,\omega^{(k)},V^{(1)},\ldots,V^{(l)})\\
&\qquad\quad\times S(\omega^{(k+1)},\ldots,\omega^{(k+m)},V^{(l+1)},\ldots,V^{(l+n)})\,.
\tag{1.58}
\end{align*}

注意，$\omega^{(i)}$和$V^{(i)}$是各不相同的对偶矢量和矢量，不是它们的分量。换句话说，先让$T$作用在适当的一组对偶矢量和矢量上，再让$S$作用在其余的上面，然后把两个答案相乘。注意，一般张量积是不可对易的：$T\otimes S\neq S\otimes T$。

现在，取基矢与基对偶矢量的张量积，就可以直截了当地构造出所有$(k,l)$张量构成的空间的基；这个基将由一切形如下式的张量组成

\begin{equation*}
\hat{e}_{(\mu_1)}\otimes\cdots\otimes\hat{e}_{(\mu_k)}
\otimes\hat{\theta}^{(\nu_1)}\otimes\cdots\otimes\hat{\theta}^{(\nu_l)}\,.
\tag{1.59}
\end{equation*}

在四维时空中，这样的基张量总共有$4^{k+l}$个。用分量记号，我们把任意的张量写成

\begin{equation*}
T=T^{\mu_1\cdots\mu_k}{}_{\nu_1\cdots\nu_l}\,
\hat{e}_{(\mu_1)}\otimes\cdots\otimes\hat{e}_{(\mu_k)}
\otimes\hat{\theta}^{(\nu_1)}\otimes\cdots\otimes\hat{\theta}^{(\nu_l)}\,.
\tag{1.60}
\end{equation*}

换个办法，我们也可以通过让张量作用在基矢和基对偶矢量上来定义分量：

\begin{equation*}
T^{\mu_1\cdots\mu_k}{}_{\nu_1\cdots\nu_l}
=T(\hat{\theta}^{(\mu_1)},\ldots,\hat{\theta}^{(\mu_k)},\hat{e}_{(\nu_1)},\ldots,\hat{e}_{(\nu_l)})\,.
\tag{1.61}
\end{equation*}

利用式(1.46)等等，你可以亲自验证这些方程都彼此相容。

与矢量一样，我们通常走捷径，用分量$T^{\mu_1\cdots\mu_k}{}_{\nu_1\cdots\nu_l}$来指代张量$T$。张量作用在一组矢量和一组对偶矢量上的方式，遵循式(1.48)确立的模式：

\begin{equation*}
T(\omega^{(1)},\ldots,\omega^{(k)},V^{(1)},\ldots,V^{(l)})
=T^{\mu_1\cdots\mu_k}{}_{\nu_1\cdots\nu_l}\,
\omega^{(1)}{}_{\mu_1}\cdots\omega^{(k)}{}_{\mu_k}\,
V^{(1)\nu_1}\cdots V^{(l)\nu_l}\,.
\tag{1.62}
\end{equation*}

于是，$(k,l)$张量有$k$个上指标和$l$个下指标。指标的次序显然是重要的，因为张量作用在它的各个宗量上的方式不必相同。

最后，张量分量在洛伦兹变换下的变换，可以借助我们已经知道的基矢和基对偶矢量的变换规律导出。答案正如你按指标位置所预期的：

\begin{equation*}
T^{\mu_1'\cdots\mu_k'}{}_{\nu_1'\cdots\nu_l'}
=\Lambda^{\mu_1'}{}_{\mu_1}\cdots\Lambda^{\mu_k'}{}_{\mu_k}\,
\Lambda^{\nu_1}{}_{\nu_1'}\cdots\Lambda^{\nu_l}{}_{\nu_l'}\,
T^{\mu_1\cdots\mu_k}{}_{\nu_1\cdots\nu_l}\,.
\tag{1.63}
\end{equation*}

这样，每个上指标都像矢量一样变换，每个下指标都像对偶矢量一样变换。

虽然我们把张量定义为从一组矢量和一组对偶矢量到$\mathbf{R}$的线性映射，但没有什么迫使我们非得作用在完整的一组宗量上。于是，$(1,1)$张量也可以作为一个从矢量到矢量的映射来作用：

\begin{equation*}
T^{\mu}{}_{\nu}:\;V^{\nu}\rightarrow T^{\mu}{}_{\nu}V^{\nu}\,.
\tag{1.64}
\end{equation*}

你可以亲自验证$T^{\mu}{}_{\nu}V^{\nu}$是一个矢量（也就是说，服从矢量的变换规律）。类似地，我们可以让一个张量作用在另一个张量上（作用在全部或部分宗量上），得到第三个张量。例如，

\begin{equation*}
U^{\mu}{}_{\nu}=T^{\mu\rho}{}_{\sigma}S^{\sigma}{}_{\rho\nu}
\tag{1.65}
\end{equation*}

就是一个完全合格的$(1,1)$张量。

考虑到这套内容的玄奥性质，你也许会担心我们对张量的介绍有些过于简略。实际上，掌握张量的概念并不需要费很大的力气；无非是把指标摆正而已，而且操作它们的规则是非常自然的。的确，有不少书喜欢把张量\emph{定义}为按式(1.63)变换的一些数的集合。虽然这在操作上是有用的，但它往往掩盖了张量作为几何实体的更深的含义——这些实体具有独立于任何所选坐标系的生命。不过，有一个我们一直含糊带过的微妙之处。对偶矢量、张量、基和线性映射这些概念属于线性代数的领地，只要手头有一个抽象的矢量空间，它们就是适用的。而在我们感兴趣的情形中，我们不只有一个矢量空间，而是时空中每一点都有一个矢量空间。我们更常感兴趣的是张量场，它可以看成时空上的取值为张量的函数。幸运的是，上面定义的那些操作，没有哪一个真正在乎我们处理的是单个矢量空间，还是每个事件一个的矢量空间的汇集。必要时，我们完全可以径直把一些东西称作$x^{\mu}$的函数，而不至于出问题。然而，你应当分清我们引入的这些概念在逻辑上的独立性，以及它们在时空和相对论中的具体应用。

在时空中，我们其实已经见过一些张量的例子，只是没有这样称呼它们。(0,2)张量最熟悉的例子就是度规$\eta_{\mu\nu}$。度规作用在两个矢量上的结果非常有用，所以它有自己的名字，叫做\textbf{内积}（inner product）（或标量积、点积）：

\begin{equation*}
\eta(V,W)=\eta_{\mu\nu}V^{\mu}W^{\nu}=V\cdot W\,.
\tag{1.66}
\end{equation*}

正如在常规的欧几里得点积中一样，我们把内积为零的两个矢量称为\textbf{正交}（orthogonal）。由于内积是标量，它在洛伦兹变换下不变；因此，任何笛卡尔惯性系的基矢——按定义选为正交——在洛伦兹变换之后仍然正交（尽管我们前面注意到它们会“剪拢”到一起）。矢量的\textbf{模}（norm）定义为矢量与自身的内积；与欧几里得空间不同，这个数不是正定的：

\begin{equation*}
\text{若 }\eta_{\mu\nu}V^{\mu}V^{\nu}\text{ 为}\;
\begin{cases}
<0\,, & \ V^{\mu}\text{ 为类时}\,,\\
=0\,, & \ V^{\mu}\text{ 为类光}\,,\\
>0\,, & \ V^{\mu}\text{ 为类空}\,.
\end{cases}
\end{equation*}

（一个矢量可以有零模长而不必是零矢量。）你会注意到，这些术语与我们前面用来分类时空中两点关系的术语相同；当然这绝非偶然，后面我们会更详细地讨论。

另一个张量是型$(1,1)$的\textbf{Kronecker $\delta$ 符号}（Kronecker delta）$\delta^{\mu}{}_{\nu}$。若把它看成从矢量到矢量（或从1-形式到1-形式）的映射，Kronecker $\delta$ 就是恒等映射（identity map）。对于这个独特的张量，我们仿照许多别的参考书的做法，把上下指标放在同一列；讲究纯粹的人也许会写成$\delta^{\mu}{}_{\rho}$或$\delta_{\rho}{}^{\mu}$，但这两者的数值完全相同，所以仅在这个例子上马虎一点不会出问题。

与Kronecker $\delta$和度规相联系的是\textbf{逆度规}（inverse metric）$\eta^{\mu\nu}$，一个型$(2,0)$的张量，（毫不奇怪地）定义为度规的“逆”：

\begin{equation*}
\eta^{\mu\nu}\eta_{\nu\rho}=\eta_{\rho\nu}\eta^{\nu\mu}=\delta^{\mu}{}_{\rho}\,.
\tag{1.67}
\end{equation*}

（它之所以是逆度规，是因为与度规相乘就得到恒等映射。）事实上，你可以验证，逆度规的分量与度规本身的分量完全相同。这只在平直空间的笛卡尔坐标中成立，在更一般的情形下就不再成立。还有\textbf{Levi-Civita符号}（Levi-Civita symbol），一个(0,4)张量：

\begin{equation*}
\tilde{\epsilon}_{\mu\nu\rho\sigma}=
\begin{cases}
+1\,, & \text{若 }\mu\nu\rho\sigma\text{ 是 0123 的偶置换}\\
-1\,, & \text{若 }\mu\nu\rho\sigma\text{ 是 0123 的奇置换}\\
0\,, & \text{其他情形}\,.
\end{cases}
\tag{1.68}
\end{equation*}

这里，“0123 的一个置换”是数0、1、2、3的一个排序，它可以由0123出发交换其中两个数字得到；偶置换通过偶数次这样的交换得到，奇置换通过奇数次得到。例如，$\tilde{\epsilon}_{0321}=-1$。（$\tilde{\epsilon}_{\mu\nu\rho\sigma}$上的波浪号，以及把它称为“符号”而不简单地称为张量，源于这样一个事实：这个对象在更一般的几何或坐标下其实并不是张量；它是一种叫做“张量密度”（tensor density）的东西。定义一个与之相关而是张量的对象是足够直截了当的，我们将把它记作$\epsilon_{\mu\nu\rho\sigma}$，并称之为“Levi-Civita张量”。讨论见第2章。）

上述张量——度规、逆度规、Kronecker $\delta$和Levi-Civita符号——的一个引人注目的性质是：尽管它们都按照张量变换规律(1.63)变换，它们的分量在平直时空的\emph{任何}惯性坐标系中都保持不变。从某种意义上说，这使得它们成为不太典型的张量例子，因为大多数张量并没有这个性质。事实上，尽管我们不去证明，具有这个性质的\emph{只有}这些张量。Kronecker $\delta$还要更加不同寻常，它在任何时空的任何坐标系中都具有完全相同的分量。从张量作为线性映射的定义来看，这很好理解；Kronecker张量可以看成从矢量到矢量（或从对偶矢量到对偶矢量）的恒等映射，显然无论坐标系如何，它都必须具有相同的分量。另一方面，度规及其逆刻画了时空的结构，而Levi-Civita符号则暗地里根本不是真正的张量。因此，当我们放弃平直时空的假设时，就必须更加小心地处理这些对象。

张量的一个更典型的例子是\textbf{电磁场强张量}（electromagnetic field strength tensor）。我们都知道，电磁场由电场矢量$E_i$和磁场矢量$B_i$构成。（记住我们用拉丁指标表示空间分量1、2、3。）实际上，这些只是在空间转动下才是“矢量”，在完整洛伦兹群下则不然。事实上，它们是一个(0,2)张量$F_{\mu\nu}$的分量，其定义为

\begin{equation*}
F_{\mu\nu}=
\begin{pmatrix}
0 & -E_1 & -E_2 & -E_3\\
E_1 & 0 & B_3 & -B_2\\
E_2 & -B_3 & 0 & B_1\\
E_3 & B_2 & -B_1 & 0
\end{pmatrix}
=-F_{\nu\mu}\,.
\tag{1.69}
\end{equation*}

从这个观点出发，通过应用式(1.63)，很容易把一个参考系中的电磁场变换成另一个参考系中的电磁场。张量形式体系的统一威力是显而易见的：与其说是两个矢量的集合，其
